% ------------------------------------------------------------
% OSTERA Template Helper Text - DO NOT EDIT THIS FILE
% This file is imported by main.tex, do not compile this file
% OTT  = OSTERA Template Text
% OTHT = OSTERA Template Help Text
% ------------------------------------------------------------

% Section 1
\newcommand{\OTHTScope}{< This section defines the intended purpose and scope of this document. The content found in this section is often reused by content providers to explain the purpose of this document in a summary form. As such, it is best to have this be shorter rather than longer. Authors and/or editors should try to keep the content at three to four paragraphs or less. If this document is intended to complement another technical specification or standards, it is often best to briefly mention that here. This section should not be used to explain the context and/or history behind the creation of this document, that type of content belongs elsewhere, either in the introduction or in an appendix. >}

% Section 4
\newcommand{\OTHTDocumentNotes}{< Add any document wide notes, disclaimers, etc. If there are no notes than one is to put "None". >}

\newcommand{\OTHTTypographicalConventions}{< Describe any typographical conventions or editorial standards that were followed when writing this document, such as fonts or highlighting that have special significance. If there are no typographical conventions than this section should simply contain "None". >}

% Section 5
\newcommand{\OTHTIntroduction}{< This section should contain information that would help the reader understand the context and other details about the document. Information about the history or what led to the creation of the document can also be included here or in an appendix. >}

\newcommand{\OTHTAnyAdditionalIntroduction}{< Any additional content that is needed or desired. >}

% Section 8
\newcommand{\OTHTSpecialConsiderations}{< This section is \textbf{REQUIRED} and \textbf{MUST} be the last numbered section in the document. This section contains information about safety, security, data protection, and privacy considerations. These can be divided up into separate subsections as desired. All documents \textbf{SHOULD} have at least a security and data protection considerations section, and otherwise \textbf{MUST} have a blank section indicating “None.” Any other considerations \textbf{MAY} also be added.

Please note that any specification that will need to register something with IANA or has plans to go on to ITU/ISO/IEC \textbf{MUST} have this section filled out. >
}

% Appendix 1
\newcommand{\OTHTAcknowledgment}{< This section should contain all of the special acknowledgements that are needed for this publication. Regardless of how many there are, the heading should be in the singular. Any individual, organization, or entity can request at any time that they not be included in this section. >}

\newcommand{\OTHTSpecialThanks}{< This section \textbf{SHOULD} include any individuals, organizations, or entities that have made material contributions to the creation and development of this document, and not just this version of the document, even if they are no longer members of the working group or association or were never members of the working group or association. >}

\newcommand{\OTHTLeadership}{< This section \textbf{SHOULD} include all of the chairs and other officers of the working group that published this document, and not just this version of the document, even if they are no longer members of the working group or the association. >}

\newcommand{\OTHTParticipants}{< This section is \textbf{OPTIONAL}, but encouraged. >}

% Appendix 2
\newcommand{\OTHTChangesFromPreviousVersion}{< The appendix \textbf{SHOULD} contain any explanatory text about the reason for this version including any major changes. The level of detail that is included in this appendix is up to the authors and/or editors. This appendix is \textbf{REQUIRED}, if there are no changes then one is to put "None." In addition to any descriptive text, all major changes \textbf{SHOULD} be in a bulleted list so that reviewers and implementers can easily understand what they need to know. >}

% End
\newcommand{\OTHTEndOfDocument}{< The following centered line represents the end of the document >}
